% ===========================================================================
\section{Emission Factors and Calculated Emissions}
\label{sec:ef}
% ===========================================================================

The spatial allocation described in Section~\ref{sec:methods} determines where emissions are
placed but does not change their national magnitude. National totals are determined only by
the activity data and the emission factors.

Table~\ref{tab:ef-full} reports the principal emission factors, their uncertainty, and their
IPCC or literature source. Table~\ref{tab:tier2} reports the country-specific and
inventory-specific parameters that replace or refine selected Tier~1 defaults. Full source
details and derivations are provided in Supplementary Note~S2.

Each factor is stored in a version-controlled data file together with its units, source and
uncertainty. This structure allows every calculated value to be traced to a documented
input.

\subsection{Emission factors}

Enteric-fermentation and manure-management factors are drawn mainly from the 2006 IPCC
defaults for the Indian Subcontinent. Managed-soils factors use the climate-specific values
in the 2019 Refinement.

The direct soil \NtwoO{} factor is the most consequential individual methodological choice.
The main inventory uses the dry-climate value $\EFone{}=0.005$. The older 2006 aggregate
value of $0.010$ and the 2019 wet-climate value of $0.016$ are retained as sensitivity
cases.

The dry-climate value is also close to the crop-weighted empirical estimate of approximately
$0.0055$ reported by \citet{cui2021} for Pakistan. That empirical value is used as a
sensitivity anchor, not as the operative emission factor.

Rice methane uses the IPCC baseline factor together with water-regime scaling. The main case
assumes continuously flooded conditions, $SF_w=1$. Intermittent irrigation and rainfed
conditions are evaluated as alternatives.

\begin{table}[!htbp]
\centering
\footnotesize
\begin{tabular}{@{}llrl@{}}
\toprule
\textbf{Factor (unit)} & \textbf{Value} & \textbf{$\pm$\%} & \textbf{Source} \\
\midrule
\multicolumn{4}{@{}l}{\emph{Enteric \CHfour{} (kg/head/yr; 2006 T10.10--10.11)}} \\
\quad Buffalo            & 55  & 40 & 2006 T10.10 \\
\quad Dairy cattle       & 58  & 50 & 2006 T10.11 \\
\quad Other cattle       & 27  & 50 & 2006 T10.11 \\
\quad Camels             & 46  & 50 & 2006 T10.10 \\
\quad Sheep, goats       & 5   & 50 & 2006 T10.10 \\
\addlinespace
\multicolumn{4}{@{}l}{\emph{Manure management (2006 T10.14--10.21)}} \\
\quad \CHfour{} buffalo, warm (kg/hd/yr) & 5 & 30 & 2006 T10.14 \\
\quad \CHfour{} dairy cattle, warm & 6 & 30 & 2006 T10.14 \\
\quad \NtwoO{} $EF_3$ solid/liquid (kg N/kg N) & 0.005 & 100 & 2006 T10.21 \\
\quad \NtwoO{} $EF_3$ drylot & 0.02 & 100 & 2006 T10.21 \\
\addlinespace
\multicolumn{4}{@{}l}{\emph{Rice \CHfour{} (2006 T5.11--5.13)}} \\
\quad Baseline $EF_c$ (kg/ha/day) & 1.30 & 53 & 2006 T5.11 \\
\quad $SF_w$ continuous (default) & 1.00 & 26 & 2006 T5.12 \\
\quad $SF_w$ intermittent (AWD) & 0.52 & 30 & 2006 T5.12 \\
\quad $SF_w$ rainfed     & 0.28 & 40 & 2006 T5.12 \\
\addlinespace
\multicolumn{4}{@{}l}{\emph{Managed soils (kg N\textsubscript{2}O-N/kg N; 2019 Refinement)}} \\
\quad \EFone{} dry (default) & 0.005 & 110 & 2019 T11.1 \\
\quad \EFone{} flooded rice  & 0.004 & 120 & 2019 T11.1 \\
\quad \EFone{} aggregate (sens.) & 0.010 & 90 & 2006 default \\
\quad \EFone{} empirical anchor & 0.0055 & 100 & \citet{cui2021} \\
\quad $EF_4$ volatilisation (dry) & 0.005 & 100 & 2019 T11.3 \\
\quad $EF_5$ leaching     & 0.011 & 100 & 2019 T11.3 \\
\addlinespace
\multicolumn{4}{@{}l}{\emph{Carbon amendments, burning}} \\
\quad Urea carbon (t\,C/t urea) & 0.20 & 10 & 2006 Eq.11.13 \\
\quad Residue burning \CHfour{} (g/kg d.m.) & 2.7 & 40 & 2006 T2.5 \\
\bottomrule
\end{tabular}
\caption{Principal emission and carbon factors, with factor-level uncertainty
($\pm$\%, 95\% relative half-width) and provenance. These factor uncertainties are
combined with activity-data uncertainty into the per-category values propagated in
the uncertainty analysis (Section~\ref{sec:methods}, \ref{app:unc}).
Values marked \emph{(sens.)} are used only in the sensitivity analysis
(Table~\ref{tab:sensitivity}); the empirical anchor adjudicates the soil-nitrogen
choice and is not the operative factor.}
\label{tab:ef-full}
\end{table}

\subsection{Country-specific and inventory-specific parameters}

The largest departures from global Tier~1 defaults occur in the livestock characterisation.

Population-average typical animal masses are used for dairy cattle, other cattle, buffalo
and camels---declared judgements resting on working pointers rather than cited
measurements, as Supplement~S2's provenance classes record (10 cited, 3 derived, 25
working pointers and 21 declared judgements across the 59 parameter rows). These values are generally higher than the corresponding IPCC defaults and affect
both nitrogen excretion and the division of enteric emissions among livestock groups. The
dairy-cattle share is derived from FAOSTAT counts of animals reported as being in milk.

Other national parameters determine how fertiliser and soil nitrogen are represented. These
include the proportion of synthetic nitrogen supplied as urea, the share applied to rice,
and the irrigated fraction used in the leaching calculation.

Table~\ref{tab:tier2} distinguishes among four types of parameter: national activity data,
Tier~2 characterisation, regional IPCC defaults, and inventory assumptions. This
category-by-category classification defines the methodological basis more accurately than
assigning one tier to the whole inventory.

\begin{table*}[!htbp]
\centering
\footnotesize
\begin{tabular}{@{}llll@{}}
\toprule
\textbf{Parameter} & \textbf{Value (unit)} & \textbf{Status} & \textbf{Source} \\
\midrule
Dairy (in-milk) cattle share & 0.295 & national activity datum & FAOSTAT QCL 2024 (16.98/57.54 M) \\
Typical mass, dairy cattle & 400 kg & Tier-2 characterisation & Sahiwal/crossbred; LRRD 2006 \\
Typical mass, other cattle & 225 kg & Tier-2 characterisation & Pak.\ adult zebu population avg. \\
Typical mass, buffalo & 380 kg & Tier-2 characterisation & Nili-Ravi herd avg.\ (LRRD) \\
Typical mass, camels  & 480 kg & Tier-2 characterisation & dromedary herd avg.; Khan 2025 \\
Nex, buffalo & 0.32 kg N/1000\,kg/d & regional IPCC default & IPCC 2006 T10.19 Asia \\
Nex, dairy cattle & 0.47 kg N/1000\,kg/d & regional IPCC default & IPCC 2006 T10.19 Asia \\
Warm-zone share (manure) & 0.70 & inventory assumption & Indus-plain climate zone \\
Irrigated cropland fraction & 0.80 & national activity datum & FAO AQUASTAT (19/22 Mha) \\
Rice share of synthetic N & 0.16 & national activity datum & rice 3.5/22 Mha cropped \\
Urea share of applied N & 0.80 & national activity datum & NFDC offtake \\
Rice cultivation period & 125 days & Tier-2 characterisation & PARC calendar (area-weighted) \\
\bottomrule
\end{tabular}
\caption{Country-specific and inventory-specific parameters that override IPCC
defaults. \emph{Status} distinguishes national activity data, genuine Tier-2
characterisation, regional IPCC defaults, and inventory assumptions --- so the
paper does not over-claim a uniform Tier-2 basis. Full provenance (source,
table/page, derivation, and uncertainty) for these and all other
parameters and factors --- management-system shares, dung-cake fraction, residue
ratios, water/amendment factors, volatilisation and leaching fractions, and the
soil \NtwoO{} factor set --- is given in Supplementary Note~S2.}
\label{tab:tier2}
\end{table*}

\subsection{Calculated emissions}

Applying the factors in Table~\ref{tab:ef-full} and the parameters in
Table~\ref{tab:tier2} to the national activity data produces the 2024 estimates reported in
Section~\ref{sec:results}.

The same emission-factor set is held constant throughout the 1990--2024 reconstruction.
Changes in the historical series therefore reflect changes in activity, including livestock
growth and increasing fertiliser use, rather than changes in the selected factors.

\ref{app:crt} presents the activity data, implied factors and emissions by gas for
the 1990 base year and the 2024 reporting year in a Common Reporting Table structure.

\subsection{Sensitivity to major methodological choices}
\label{sec:sensitivity}

Three assumptions account for most of the methodological sensitivity: the direct soil factor
$\EFone{}$, the rice water-regime factor $SF_w$, and the distribution of livestock manure
among management systems.

Increasing $\EFone{}$ from the main-case value of $0.005$ to the 2006 aggregate value of
$0.010$ raises the national total by \SI{11.9}{}~Mt~\COe{}. Using the empirical value of
$0.0055$ produces little change and gives a managed-soils estimate of \SI{29.7}{}~Mt, close
to that of \citet{cui2021}.

Changing the rice water regime from continuously flooded to intermittent irrigation reduces
the national total by \SI{7.9}{}~Mt~\COe{}. A fully rainfed assumption reduces it by
\SI{11.9}{}~Mt~\COe{}.

Using the GLEAM~3.0 South-Asia manure-management distribution \citep{gleam2022} raises the
national total by \SI{9.9}{}~Mt~\COe{}. Unlike the soil and rice cases, this scenario
requires a complete model run because manure management is coupled to the quantity of manure
nitrogen later applied to soils. It therefore changes both manure-management and
managed-soils emissions.

All tested scenarios remain within the main Monte-Carlo interval and do not change the
ranking of the leading emission categories.

\begin{table}[!htbp]
\centering
\footnotesize
\begin{tabular}{@{}lccc@{}}
\toprule
\textbf{Scenario} & \textbf{Rice} & \textbf{Soils} & \textbf{Total} \\
\midrule
Main case ($\EFone{}=0.005$; rice flooded) & 16.5 & 28.5 & 223.5 \\
\quad $\EFone{}=0.010$ (2006 aggregate) & 16.5 & 40.4 & 235.4 \\
\quad $\EFone{}=0.016$ (2019 wet)       & 16.5 & 54.7 & 249.7 \\
\quad $\EFone{}=0.0055$ (Cui 2021)      & 16.5 & 29.7 & 224.6 \\
\quad $EF_3PRP$ wet (2019)$^{b}$        & 16.5 & 30.3 & 225.3 \\
\quad $EF_3PRP$ aggregate (2006)        & 16.5 & 42.8 & 237.7 \\
\quad Both 2006 direct factors          & 16.5 & 54.7 & 249.6 \\
\quad Rice intermittent $SF_w=0.52$ (AWD) &  8.6 & 28.5 & 215.5 \\
\quad Rice rainfed $SF_w=0.28$          &  4.6 & 28.5 & 211.5 \\
\quad Manure GLEAM 3.0 split$^{a}$      & 16.5 & 27.6 & 233.4 \\
\bottomrule
\end{tabular}
\caption{Sensitivity of the 2024 national agricultural total (Mt~\COe, AR5) to the
principal methodological choices. $^{a}$Soil and rice alternatives can be recalculated directly because each touches a
single category and the affected terms scale linearly. The GLEAM manure-management case
requires a complete model run: the nitrogen ledger is coupled, so a different split changes how much manure nitrogen
reaches the field and moves managed soils (\SI{28.5}{} to \SI{27.6}{}~Mt) as well as manure
management (\SI{21.2}{} to \SI{32.1}{}~Mt). It is the same scenario as
Table~\ref{tab:apstruct} and reports the same total. $^{b}$The pasture factor had no
scenario beside it before this release, which was an asymmetry rather than a judgement: it is
the larger lever of the two soils factors---the 2019 Refinement cut the cattle value tenfold
from the 2006 default where it halved \EFone{}---and it carries two-thirds of the difference
against FAOSTAT's managed soils. The two factor scenarios are disjoint, so the combined 2006
case is their sum rather than a separate run. Raised via the GLEAM 3.0 South-Asia management-system split (a hidden
uncertainty surfaced by the cross-check, not adopted). Generated by
\texttt{src/pkinv/eval/sensitivity\_table.py}.}
\label{tab:sensitivity}
\end{table}
